Considera les rectes $r\colon\dfrac{x+1}{1}=\dfrac{y-2}{-1}=\dfrac{z-1}{2}$ i
$s\colon\left\{\begin{aligned}x&=-1+2\mu\\y&=-1+\mu\\z&=6-\mu\end{aligned}\right.$

\begin{apartats}

\apartat[1,25]{0,75}
Comprova que les rectes es tallen i troba'n el punt de tall.

\begin{solucio}
$\vec u=(1,-1,2)$ i $\vec v=(2,1,-1)$ no són proporcionals. En paramètriques, $r$ és
$(-1+t,\ 2-t,\ 1+2t)$. Igualant: $-1+t=-1+2\mu$, $2-t=-1+\mu$ i $1+2t=6-\mu$. Les dues primeres donen
$t=2\mu$ i $2-2\mu=-1+\mu$, és a dir, $\mu=1$ i $t=2$, que compleixen la tercera ($5=5$). Es tallen al
punt $(1,0,5)$.
\end{solucio}

\begin{tria}{rectes-secants}
\itemtria{pla-que-les-conte}{1}{1,25}
Troba l'equació del pla que conté les dues rectes.

\begin{solucio}
El pla passa per $P(-1,2,1)$, de $r$, i té com a vectors directors $\vec u$ i $\vec v$:
\[
\begin{vmatrix}x+1&y-2&z-1\\1&-1&2\\2&1&-1\end{vmatrix}=-(x+1)+5(y-2)+3(z-1)=0
\ \Longrightarrow\ x-5y-3z+14=0.
\]
Comprovació: el punt $(-1,-1,6)$ de $s$ dona $-1+5-18+14=0$.
\end{solucio}

\itemtria{amb-la-diagonal}{1}{1,25}
Estudia la posició relativa de la recta $r$ i la recta $x=y=z$.

\begin{solucio}
La recta $x=y=z$ passa per l'origen $O$ amb vector director $(1,1,1)$, que no és proporcional a
$\vec u=(1,-1,2)$. Amb $\overrightarrow{PO}=(1,-2,-1)$, on $P(-1,2,1)$ és de $r$:
\[
\begin{vmatrix}1&-1&2\\1&1&1\\1&-2&-1\end{vmatrix}=-7\ne0.
\]
El rang és 3: les dues rectes es creuen.
\end{solucio}
\end{tria}

\begin{nomesllarg}
\apartat{0,75}
Escriu les equacions d'una recta que es creui amb $r$. Justifica-ho.

\begin{solucio}
Resposta oberta. Per exemple, l'eix $OX$: passa per l'origen amb vector director $(1,0,0)$, que no és
proporcional a $\vec u$. Amb $\overrightarrow{PO}=(1,-2,-1)$:
$\begin{vmatrix}1&-1&2\\1&0&0\\1&-2&-1\end{vmatrix}=-5\ne0$. El rang és 3: l'eix $OX$ i $r$ es creuen.
\end{solucio}
\end{nomesllarg}

\end{apartats}
